% =============================================================================
% chapters/02_literature.tex — Literature Review
% =============================================================================

\section{Literature Review}

This chapter provides a comprehensive review of the theoretical foundations and empirical evidence relevant to your study. The review should be organized into thematic areas that align with your research questions, identifying what is already known and where the gaps lie.

\subsection{Theoretical Foundations}
Introduce the main theories or frameworks underpinning your research. For example, in management research, scholars frequently draw upon the dynamic capabilities framework proposed by \textcite{Teece2007} or the resource-based view \parencite{Barney1991}.

\subsection{Key Concepts}
Discuss the core concepts of your study. For example, if you are examining digital transformation, you can reference systematic reviews such as \textcite{Vial2019} to define the processual nature of digital shifts and distinguish them from mere digitization \parencite{Verhoef2021}.

\subsection{Empirical Review}
Review previous empirical studies that are directly relevant to your topic. Focus on their methodologies, key findings, and the limitations they identified.

\subsection{Research Gaps and Contribution}
Conclude the literature review by explicitly highlighting the research gaps that your study aims to fill. Describe how your work contributes to the academic and practical understanding of this domain.
